\chapter{Project management command}

%---------------------------------------------
%---------------------------------------------
%---------------------------------------------
\section{Readible file formats}

\subsection{Readible/Binary}

As photogrammetric pipeline is a complex process, made of several computation,
the result of each computation has to be writen in some files that
will be read at by next process. There is basically two familly of such files:

\begin{itemize}
   \item files that may have some interest to  be read or manipulated by human
         or other programm, in this case {\tt MMVII} standar tagged format
	 as {\tt xml} of {\tt json}; example of such file are calibration
	 or pose estimation;

   \item files that have probably low interest for human, and for efficiency 
         they are stored in binary format  (note that there exist also a
         text version of this binary format, note easy to read, but that facilitate
	   import export);
\end{itemize}

\subsection{Xml and Json files}

{\tt MMVII} offer the possibility to export data in two different tagged format : {\tt xml} and
{\tt json}. By default it's {\tt xml}, but this can be changed using mecanism descrined in \ref{UserParametrisation}.
Implicit conversion may appear in a near future for sharing files between user,
by the way it's recommanded that user make a choice once for all  to avoid any problem.

The folder {\tt MMVII-UseCaseDataSet/SampleFiles} contains examples
of files in {\tt xml} and {\tt json}. 

Note that {\tt xml} allows real coments and, for example, in the file {\tt Calib...xml}, they are used to
explicit  the meanin of distorsion parameters. As {\tt json} do not have
real comments, we use the special tags , for example in file
{\tt Calib...json} you can find {\tt "<!--comment6-->":"(X,0)"}
(corresponding to  {\tt <!--(X,0)-->} in {\tt xml}).

These file are made to be relatively easy to read. They can also be created or
modified easily by user or programm however there are some \emph{strict} rules to observe so that such
file remain valid . Starting from a valid file, here are example of things that can be done to maintain
validity :

\begin{itemize}
       \item change the value of atoms while respecting their type (float, int, string \dots);

       \item add or supress an element  in a sequence or a map,  see for example  the file
	       {\tt TestObj.*}, the element of sequence are tagged {\tt el} in {\tt xml}, while the element of
		pair key-value  in a map are tagged {\tt K/V};

       \item supress or add any comment;

       \item supress or add an optionnal value (rare case for now), the optionnal value can
	     be  detected as their tags begin by {\tt Opt:}, see file {\tt F\_T2.*};
\end{itemize}

And here a \emph{non-exhautive} list of thing you cant do :

\begin{itemize}
        \item obviously create a non-valide xml/json file;

	\item swap two elements of different tags  \emph{very bad, dont do it, strictly forbiden, naughty , Micmac is whatching you \dots}
          also it woul be theoretically possible to recover the information, this would involve unnecessary software devlopment
	  that wont be done;

        \item supress a non optional tag ;

	\item add/supress a value in fixed size tab (used to represent point for exeample);
\end{itemize}

The format of each file is :

\begin{itemize}
	\item a single root node , named {\tt root} in {\tt xml} and anonymous in Jason,
	\item root node has exacly $3$ sub-node that have  a fixed tag

	\item the first node is the type, it must be tagged {\tt Type} and its value must {\tt MMVII\_Serialization} 
              signing the fact that is was created by {\tt MMVII};

      \item the second node is a version number, it must be tagged {\tt Version}, its value will allow
	    compatibility policy with older files, unused at the time being;
             
    \item the third node must be tagged {\tt Data} and contains in fact the data itself !
          in most frequent case, it will contain a single node (see {\tt Calib*, Ori*} ) allowing
          some type-checking at the very begining   of the command (i.e we can check that {\tt Calib*}
          are most probably  MMVII-calibration files as their data contains the single tag {\tt InternalCalibration});
\end{itemize}


%---------------------------------------------
%---------------------------------------------
%---------------------------------------------


\section{User specific parametrization}

\label{UserParametrisation}

Sometimes a user needs to fix stable default parametrization of {\tt MMVII}
without repeating the same options in every command line. This mechanism is
called a \emph{user profile}. A profile is a set of key/value preferences stored
outside the project data directory and loaded by every {\tt MMVII} command.

The parameters accessible for now are :

\begin{itemize}
    \item maximal number of processors allowed when {\tt MMVII} executes
          parallel computation, through the key {\tt NbProcMax};

    \item default format for serialized tagged data, through the key
          {\tt TaggedSerialMode};

    \item default format for serialized vector-like data, through the key
          {\tt VectSerialMode};

    \item command used by {\tt MMVII} to open generated {\tt pdf} files, through
          the key {\tt PdfOpen}.
\end{itemize}

The distributed defaults are stored in the installation directory, in
{\tt MMVII-LocalParameters/MMVII-Default-Profile.xml}. This file is shared with
{\tt MMVII} and should be modified only when changing the distributed defaults.
For a standard installation it contains a map of string key/value pairs such as :

\begin{verbatim}
<?xml version="1.0" encoding="UTF-8" standalone="yes" ?>
<Root>
   <Type>"MMVII_Serialization"</Type>
   <Version>"0.0.0"</Version>
   <Data>
      <Pair>
         <K>"NbProcMax"</K>
         <V>"4"</V>
      </Pair>
      <Pair>
         <K>"TaggedSerialMode"</K>
         <V>"xml"</V>
      </Pair>
      <Pair>
         <K>"VectSerialMode"</K>
         <V>"csv"</V>
      </Pair>
      <Pair>
         <K>"PdfOpen"</K>
         <V>""</V>
      </Pair>
   </Data>
</Root>
\end{verbatim}

The user-specific files are not stored in the project folder. They are stored in
an operating-system-specific configuration directory :

\begin{itemize}
    \item on Linux, {\tt \$XDG\_CONFIG\_HOME/MMVII} when
          {\tt XDG\_CONFIG\_HOME} is an absolute path, otherwise
          {\tt \$HOME/.config/MMVII};
    \item on macOS, {\tt \$HOME/Library/Application Support/MMVII};
    \item on Windows, {\tt \%APPDATA\%/MMVII}.
\end{itemize}

This directory is also used by {\tt MMVII} for other user-local data, including the user camera database. For profiles, it contains two kinds of files. The file
{\tt MMVII-Current-Profile.xml} stores the name of the current profile. Each
profile is then stored in a file named {\tt MMVII-profile-<Name>.xml}. Profile
names must contain only letters, digits, underscores or hyphens, and must not be
empty. If {\tt MMVII-Current-Profile.xml} does not exist, {\tt MMVII} uses the
profile named {\tt Default}.

When a profile is loaded, {\tt MMVII} first reads
{\tt MMVII-Default-Profile.xml}, then overlays the values found in
{\tt MMVII-profile-<Name>.xml} if this user file already exists. The merged
profile is written back to the user profile file. This makes new distributed keys
appear automatically in existing user profiles, while preserving user overrides.

The command {\tt EditProfile} is the normal way to inspect and modify profiles.
Called without option, it prints the known profile names, the current profile
between square brackets, the current profile file path, and the key/value pairs. This path also shows the user configuration directory :

\begin{verbatim}
MMVII EditProfile
\end{verbatim}

To make a profile current for subsequent commands, use {\tt SetCurrent} :

\begin{verbatim}
MMVII EditProfile SetCurrent=Default
MMVII EditProfile SetCurrent=Test
\end{verbatim}

To set one key/value pair in the selected profile, use {\tt KeyVal}. It must
contain exactly two values, the key and the value. To specify an empty string for
a value, use the syntax {\tt KeyVal=[Key,,]}.

\begin{verbatim}
MMVII EditProfile KeyVal=[NbProcMax,8]
MMVII EditProfile KeyVal=[TaggedSerialMode,json]
MMVII EditProfile KeyVal=[VectSerialMode,csv]
MMVII EditProfile KeyVal=[PdfOpen,evince]
MMVII EditProfile KeyVal=[Empty,,]
\end{verbatim}

To remove a key from the selected profile, use {\tt DelKey} :

\begin{verbatim}
MMVII EditProfile DelKey=PdfOpen
\end{verbatim}

Removing a key from a user profile does not remove it from the distributed
default profile. If the key exists in {\tt MMVII-Default-Profile.xml}, it will be
restored with its default value at the next synchronization unless the user
profile overrides it again.

Every {\tt MMVII} command also has a global option {\tt Profile}. It selects a
profile only for the current execution and does not modify
{\tt MMVII-Current-Profile.xml} :

\begin{verbatim}
MMVII SomeCommand ... Profile=Test
\end{verbatim}

Command-line options remain authoritative for the command being executed. For
example, {\tt NbProcMax} is applied only when the command did not explicitly
initialize the global option {\tt NbProc}. When both are present, the explicit
command-line option keeps priority.


%---------------------------------------------
%---------------------------------------------
%---------------------------------------------

\section{General data organization}

The data organization in {\tt MMVII} is the following :

\begin{itemize}
     \item for a given project, almost all the file created are located somewhere under
           the same folder {\tt MMVII-PhgrProj},  this is necessary because during a complex 
           photogrammetric process many files will be created and we dont want to encumber
           the main folder;

     \item an example of such folder (resulting from the coded-target usecase) can be
	     found under {\tt MMVII-UseCaseDataSet/SampleFiles}

     \item for each kind of processing, there exist a subfolder corresponding to the "nature"
            of the data stored;

    \item  in the example there is of folder for orientation {\tt Ori}, one for points 3d coordinates
    {\tt ObjCoordWorld}, one for point measurement
	   {\tt ObjMesInstr}, one fore handling meta data {\tt MetaData}, one for storing reports
           {\tt Reports}, all the name of this subfolder are defined by {\tt MMVII} and cannot be changed
            by the user;

    \item there is (will be) many other king of folder : homologous point, radiometric model, radiometric data,
           \dots

     \item in each of the predefined folder, there exist different subfolder corresponding to different step
           of the process; the name of this subfolder are specified by user, sometime as input to a command,
	   sometime as output to a command; typically the output a command being the input of the next command;
\end{itemize}

In this example, we have $4$ different folder for orientation in {\tt Ori} :

\begin{itemize}
      \item  {\tt 11P} which is the result of initial estimation using uncalibrated spaced resection
	      (with the "$11$ parameters method");

      \item  {\tt Resec} which is the result of pose estimation using calibrated camera, using as input
	      the calibration stored in {\tt 11P};

      \item  {\tt BA} which is the result of pose estimation using bundled adjusment, using as input
	      for initial value the pose stored in {\tt Resec};

      \item   the refined pose of  {\tt BA} are used as input to drive a research of uncoded target with
	      accurate initial position in imahe;
 
       \item using the additional target, {\tt BA} is used as input to a new bundle adjusment and the result is
             stored in {\tt BA2}.

\end{itemize}


%---------------------------------------------
%---------------------------------------------
%---------------------------------------------


\section{Help command}
\index{Help}

\label{HelpCmd}


\section{Bench command}


\begin{itemize}
    \item {\tt  MMVII Bench 2 }  : standard mode  for execuring all bench at level 2

    \item {\tt MMVII Bench 2 PatBench=.*Der.* Show=0}  : execute benches matchin {\tt ".*Der.*"},
          {\tt  Show} is explicite as, by default, it is set to {\tt  true} 
          when {\tt  PatBench} is set;
   
    \item {\tt MMVII Bench 1 PatBench=XXX} : pattern specified but no match, print all benche existing


    \item {\tt MMVII Bench 2 KeyBug=Debord\_M1 }  : force the generation of  a given error

    \item {\tt MMVII Bench 1 KeyBug=XXXX }  : will print all possible value for explicit error generation


    \item {\tt MMVII Bench 1 PatBench=InspectCube }  : as InspectCube is not a bench function, but 
         only print information, exact name must be set with {\tt PatBench}

\end{itemize}






{\tt MMVII Bench 2 KeyBug=XXX }  : standard mode  for execuring all bench at level 2

{\tt MMVII Bench 1 PatBench=MemoryOperation KeyBug} : pattern specified but no match, print all benche existing


  %---------------------------------------------
%---------------------------------------------
%---------------------------------------------

\section{Image metadata calculator : \texttt{EditCalcMTDI}}
\index{EditCalcMTDI}

\subsection{Purpose}

When images lack embedded metadata (EXIF information), {\tt MMVII} needs the user
to supply camera parameters such as focal length, aperture, or camera model.
The command {\tt EditCalcMTDI} lets the user define and edit a set of
\emph{pattern-based rules} stored in a project file; these rules are then applied
automatically to associate each image filename with the correct metadata value.

Rules are written in the file {\tt CalcMTD.xml} located in
{\tt MMVII-PhgrProj/MetaData/<DirIn>/}.
Each rule is a pair (regular expression pattern, substitution value).
When a metadata attribute is queried for a given image filename, the rules are
tested in order and the first matching rule wins.

\subsection{Syntax}

\begin{verbatim}
	== Mandatory unnamed args : ==
	* string  :: MetaData input directory (e.g. Std)
	* eMTDIm  :: Type of metadata attribute to edit
	{Focalmm, FocalPix, PPPix, Aperture, ModelCam, NbPixel, AdditionalName}

	== Optional named args : ==
	* [Name=ImTest]        string              :: Image filename to test the rules against
	* [Name=Show]          bool                :: Print all existing rules for the chosen type
	* [Name=Save]          bool                :: Save modifications to CalcMTD.xml
	* [Name=Modif]         [string,string,int] :: Rule to add/replace: [Pattern, Substitution, Rank]
\end{verbatim}

\subsection{Metadata attribute types}

The second mandatory argument selects which metadata attribute to edit:

\begin{itemize}
	\item {\tt Focalmm} --- focal length in millimetres;

	\item {\tt FocalPix} --- focal length expressed directly in pixels;

	\item {\tt PPPix} --- principal point position in pixels (2D point);

	\item {\tt Aperture} --- aperture value (f-stop), used by the radiometric model;

	\item {\tt ModelCam} --- camera model name, e.g.\ {\tt "NIKON D5600"}, used to query the
	internal camera database;

	\item {\tt NbPixel} --- image dimensions in pixels, useful when no camera database entry
	is available (e.g.\ after rescaling);

	\item {\tt AdditionalName} --- an extra string to distinguish cameras that share the same
	model name and focal length. This is typically needed when the dataset contains a
	\emph{block of cameras} (e.g.\ a multi-camera rig or a synchronised aerial system)
	where several camera bodies are identical in model and focal length but must receive
	separate calibrations. The {\tt AdditionalName} value groups images by camera body
	so that each body is treated as a distinct sensor.
\end{itemize}

\subsection{The \texttt{Modif} argument}

{\tt Modif} takes exactly three values: {\tt [Pattern, Substitution, Rank]}.

\begin{itemize}
	\item {\tt Pattern} is a C++ regular expression matched against the image filename;

	\item {\tt Substitution} is the value to assign when the pattern matches.
	It may contain C++ regex back-references (e.g.\ {\tt \$1} for the first
	captured group);

	\item {\tt Rank} controls where the rule is inserted in the ordered list:
	\begin{itemize}
		\item $\geq 0$ : insert or replace at position {\tt Rank} (0-based); if {\tt Rank}
		is beyond the current list length the rule is appended;
		\item $-1$ : prepend the rule (highest priority);
		\item $-2$ : replace the entire rule list with this single rule.
	\end{itemize}
\end{itemize}

Modifications are only persisted when {\tt Save=1} is passed.  Before overwriting,
{\tt MMVII} keeps up to 5 backup copies of the previous {\tt CalcMTD.xml}.

\subsection{Testing rules with \texttt{ImTest}}

{\tt ImTest} is optional: the command runs without it and is useful on its own
to add a rule ({\tt Modif=\dots~Save=1}) or inspect the current rule set ({\tt Show=1}).
When {\tt ImTest=<filename>} is supplied, {\tt MMVII} additionally runs the full
matching procedure on that filename and prints a trace showing each rule that was
tested and whether it matched, finishing with the final computed value.
This makes it easy to verify that a newly added rule behaves as expected before
saving, or to diagnose why a particular image is receiving an unexpected value.

\subsection{Examples}

Assign the camera model {\tt "NIKON D5600"} to every {\tt .JPG} file and verify
the result on a test image:

\begin{lstlisting}
	MMVII EditCalcMTDI Std ModelCam ImTest=043_0136.JPG \
	Modif=[.*.JPG,"NIKON D5600",0] Save=1
\end{lstlisting}

Assign a focal length of 24\,mm to rescaled {\tt .tif} files:

\begin{lstlisting}
	MMVII EditCalcMTDI Std Focalmm ImTest=043_0005_Scaled.tif \
	Modif=[.*_Scaled.tif,24,0] Save=1
\end{lstlisting}

Use a regex capture group to derive the {\tt AdditionalName} from the filename
(everything before the first underscore):

\begin{lstlisting}
	MMVII EditCalcMTDI Std AdditionalName ImTest=043_0005_Scaled.tif \
	Modif=["(.*)_.*","$1",0] Save=1
\end{lstlisting}

Display all rules currently defined for focal length, without modifying anything:

\begin{lstlisting}
	MMVII EditCalcMTDI Std Focalmm Show=1
\end{lstlisting}

